%%TPotEqBolton
% Calcul de la temp�rature potentielle �quivalente selon la formule de BOLTON
% � partir de Ps, Ts et Hum 

Conversion �ventuelle de la temp�rature $T_s$ en Kelvin \newline
Calcul du rapport de m�lange $r_{mel}$ = \texttt{rmel\_hu($Hum$, $T_s$, $P_s$)} (en g/kg)\newline
Calcul de la temp�rature potentielle $\theta$ = \texttt{tpot\_{Bolton}($P_s$, $T_s$, $Hum$)} \newline
Calcul de la temp�rature du point de condensation $T_{cond}$ = \texttt{tcond($T_s$, $Hum$)} 
\[
Tpot_{eq} = \theta . \exp \left( \left( \frac{3.376}{T_{cond}} - 0.00254 \right) . r_{mel}. \left( 1 + 0.00081 . r_{mel} \right) \right) 
\]
Si la temp�rature $T_s$ �tait en degr�s Celsius, conversion de $Tpot_{eq}$ en degr�s Celsius
